\documentclass[11pt,a4paper]{article}
\usepackage[margin=2.5cm]{geometry}
\usepackage{amsmath,amssymb,amsthm}
\usepackage{stmaryrd}
\usepackage{graphicx}
\usepackage{hyperref}
\usepackage{enumitem}
\usepackage{booktabs}

\newcommand{\Set}{\mathbf{Set}}
\newcommand{\Cat}{\mathbf{Cat}}
\newcommand{\Cof}{\mathbf{Cof}}
\newcommand{\DCont}{\mathbf{DCont}}
\newcommand{\cont}[2]{#1 \triangleleft #2}
\newcommand{\id}{\mathrm{id}}
\newcommand{\cod}{\mathrm{cod}}
\newcommand{\dom}{\mathrm{dom}}
\newcommand{\ob}{\mathrm{ob}\,}
\newcommand{\sharpp}{\sharp}
\newcommand{\co}{\mathrel{\text{\reflectbox{$\rightharpoonup$}}}}
\DeclareMathOperator{\Hom}{Hom}

\theoremstyle{definition}\newtheorem{definition}{Definition}[section]
\theoremstyle{plain}\newtheorem{proposition}[definition]{Proposition}
\newtheorem{theorem}[definition]{Theorem}
\newtheorem{lemma}[definition]{Lemma}
\newtheorem{corollary}[definition]{Corollary}
\theoremstyle{remark}\newtheorem{example}[definition]{Example}
\newtheorem{remark}[definition]{Remark}

\title{\textbf{The Morphisms of Directed Containers are Cofunctors:\\
completing the dictionary $\DCont\cong\Cof$}}
\author{MacBeth}
\date{9 June 2026}

\begin{document}
\maketitle

\begin{abstract}
The object-level dictionary of \texttt{directed-categories.tex} identifies
directed containers with small categories. We extend it to morphisms. The
container morphism $(f,f^\sharp)$ has its position component
$f^\sharp_s\colon P'(fs)\to P(s)$ \emph{contravariant} on positions, whereas a
functor acts \emph{covariantly} on morphisms; the two variances are opposite,
so directed-container morphisms cannot be functors. We resolve the mismatch:
with the directed-structure compatibility conditions made fully explicit,
$(f,f^\sharp)$ is exactly a \emph{cofunctor} (retrofunctor) of the associated
categories. The assignment is a bijection on hom-sets and a functor, and it is
bijective on objects, so it is an \emph{isomorphism of categories}
$\DCont\cong\Cof$, where $\Cof$ is the category of small categories and
cofunctors. The would-be statement $\DCont\simeq\Cat$ is therefore false: the
object dictionary does not extend to a full functor into $\Cat$, since
$\Cof$ and $\Cat$ have genuinely different hom-sets over their common objects
(neither contains the other). As a foil we show that the \emph{opposite}
variance --- a covariant position map $P(s)\to P'(fs)$ --- recovers functors
exactly. All counts are verified combinatorially across $36$ ordered pairs of
small categories.
\end{abstract}

\section{Setup and the variance problem}

We adopt the notation of \texttt{directed-categories.tex}. A directed container
$D=(\cont SP,\mathsf o,\downarrow,\oplus)$ determines a small category
$\mathcal C(D)$ by the dictionary
\[
  \ob\mathcal C(D)=S,\qquad
  \mathcal C(D)(s,s')=\{p\in P(s): s\downarrow p=s'\},\qquad
  \id_s=\mathsf o(s),\qquad q\circ p = p\oplus q,
\]
so that $P(s)=\coprod_{s'}\mathcal C(D)(s,s')$ is the set of all morphisms
\emph{out of} $s$, and $s\downarrow p=\cod(p)$. Conversely every small category
$\mathcal C$ arises as $\mathcal C(D)$ for a unique directed container $D$;
this object-level bijection is Theorem~4.1 of \texttt{directed-categories.tex}.

A morphism of the underlying containers
$(\cont SP)\to(\cont{S'}{P'})$ is a pair
\[
  (f,f^\sharp),\qquad f\colon S\to S',\qquad
  f^\sharp_s\colon P'(fs)\to P(s)\ \ (s\in S),
\]
the data of a natural transformation between extension functors. Note the
direction: $f$ goes forward on shapes, while $f^\sharp$ pulls positions
\emph{backward}, from the target shape $fs$ to the source shape $s$. Under the
dictionary, $f^\sharp_s$ sends a morphism out of $fs$ in $\mathcal C(D')$ to a
morphism out of $s$ in $\mathcal C(D)$.

\paragraph{The problem.} A functor $F\colon\mathcal C\to\mathcal C'$ is
\emph{covariant} on hom-sets: it sends a morphism out of $s$ to a morphism out
of $Fs$, a map $P(s)\to P'(Fs)$. The container component $f^\sharp$ runs the
other way, $P'(fs)\to P(s)$. The variances are opposite. Hence the naive hope
``a $\DCont$-morphism is a functor'' is ill-typed, and the assertion
$\DCont\simeq\Cat$ must be examined rather than assumed. We show the correct
target is the category of cofunctors.

\section{Cofunctors}

\begin{definition}[Cofunctor]\label{def:cof}
A \emph{cofunctor} $\varphi\colon\mathcal C\co\mathcal D$ between small
categories consists of
\begin{itemize}[nosep]
  \item an object map $\varphi_0\colon\ob\mathcal C\to\ob\mathcal D$;
  \item a \emph{lift}: for every object $c$ of $\mathcal C$ and every morphism
        $h\colon\varphi_0 c\to d$ of $\mathcal D$, a morphism
        $\varphi^\uparrow(c,h)\colon c\to c'$ of $\mathcal C$;
\end{itemize}
subject to:
\begin{align}
  &\text{(C0)} && \varphi_0\big(\cod\varphi^\uparrow(c,h)\big)=\cod(h);
        &&\text{(codomain lift)}\notag\\
  &\text{(C1)} && \varphi^\uparrow(c,\id_{\varphi_0 c})=\id_c;
        &&\text{(unit)}\notag\\
  &\text{(C2)} && \varphi^\uparrow(c,h'\circ h)
                   =\varphi^\uparrow(c',h')\circ\varphi^\uparrow(c,h),
        \quad c'=\cod\varphi^\uparrow(c,h). &&\text{(composition)}\notag
\end{align}
In (C2), $h\colon\varphi_0 c\to d$ and $h'\colon d\to d''$; the term
$\varphi^\uparrow(c',h')$ typechecks because $\varphi_0 c'=d=\dom(h')$ by (C0).
Cofunctors compose --- object maps forward, lifts back ---
$(\psi\varphi)^\uparrow(c,m)=\varphi^\uparrow(c,\psi^\uparrow(\varphi_0 c,m))$
--- with identities $(\id_0,\id^\uparrow)$, forming a category $\Cof$ of small
categories and cofunctors.
\end{definition}

Cofunctors are the comonoid morphisms in $(\mathbf{Poly},\triangleleft)$
(Ahman--Uustalu; Garner); we revisit this in \S\ref{sec:poly}.

\section{The morphism dictionary}

We now write the compatibility conditions that promote a container morphism
$(f,f^\sharp)$ to a morphism of \emph{directed} containers, and read them as
the cofunctor laws. For $h\in P'(fs)$, recall $f^\sharp_s(h)\in P(s)$ and
$(fs)\downarrow' h\in S'$.

\begin{definition}[Directed-container morphism]\label{def:dcm}
A \emph{morphism of directed containers} $D\to D'$ is a container morphism
$(f,f^\sharp)$ satisfying, for all $s\in S$, $h\in P'(fs)$ and
$k\in P'\big((fs)\downarrow' h\big)$:
\begin{align}
  &\text{(M0)} && f\big(s\downarrow f^\sharp_s(h)\big) = (fs)\downarrow' h;
        &&\text{($\downarrow$-compatibility)}\notag\\
  &\text{(M1)} && f^\sharp_s\big(\mathsf o'(fs)\big) = \mathsf o(s);
        &&\text{(root preservation)}\notag\\
  &\text{(M2)} && f^\sharp_s(h\oplus' k)
        = f^\sharp_s(h)\ \oplus\ f^\sharp_{\,s\downarrow f^\sharp_s(h)}(k).
        &&\text{(shift homomorphism)}\notag
\end{align}
\end{definition}

\begin{remark}[Well-typedness of (M2)]
The left side: $h\oplus' k\in P'(fs)$ since $\oplus'$ lands in $P'(fs)$, so
$f^\sharp_s$ applies. The right side: put $p:=f^\sharp_s(h)\in P(s)$. By (M0),
$(fs)\downarrow' h = f(s\downarrow p)$, so $k\in P'(f(s\downarrow p))$ and
$f^\sharp_{s\downarrow p}(k)\in P(s\downarrow p)$; then
$p\oplus f^\sharp_{s\downarrow p}(k)\in P(s)$. Thus (M2) typechecks precisely
because of (M0) --- mirroring how D5 typechecks because of D4.
\end{remark}

\begin{lemma}[Hom-set bijection]\label{lem:hom}
Fix directed containers $D,D'$ with categories $\mathcal C=\mathcal C(D)$,
$\mathcal C'=\mathcal C(D')$. The map
\[
  (f,f^\sharp)\ \longmapsto\ \big(\varphi_0:=f,\ \varphi^\uparrow(s,h):=f^\sharp_s(h)\big)
\]
is a bijection between morphisms $D\to D'$ of directed containers
(Definition~\ref{def:dcm}) and cofunctors
$\mathcal C\co\mathcal C'$ (Definition~\ref{def:cof}).
\end{lemma}

\begin{proof}
Under the dictionary, a position $h\in P'(fs)$ \emph{is} a morphism out of $fs$
in $\mathcal C'$, with codomain $(fs)\downarrow' h$; and $f^\sharp_s(h)\in P(s)$
\emph{is} a morphism out of $s$ in $\mathcal C$, with codomain
$s\downarrow f^\sharp_s(h)$. So the data $(f,f^\sharp)$ and $(\varphi_0,
\varphi^\uparrow)$ are literally the same data. We match the conditions
term-by-term.

\emph{(M0)$\Leftrightarrow$(C0).} (M0) reads
$f(s\downarrow f^\sharp_s(h))=(fs)\downarrow' h$. The left side is
$\varphi_0(\cod\varphi^\uparrow(s,h))$ and the right side is $\cod(h)$ in
$\mathcal C'$. This is exactly (C0).

\emph{(M1)$\Leftrightarrow$(C1).} The identity at $fs$ is
$\mathsf o'(fs)=\id_{fs}$, and the identity at $s$ is $\mathsf o(s)=\id_s$. So
$f^\sharp_s(\mathsf o'(fs))=\mathsf o(s)$ is
$\varphi^\uparrow(s,\id_{\varphi_0 s})=\id_s$, exactly (C1).

\emph{(M2)$\Leftrightarrow$(C2).} Translate $\oplus$ into $\circ$ by
$x\oplus y=y\circ x$. With $p=f^\sharp_s(h)$ and $s'=s\downarrow p$, the right
side of (M2) is $p\oplus f^\sharp_{s'}(k)=f^\sharp_{s'}(k)\circ f^\sharp_s(h)
=\varphi^\uparrow(s',k)\circ\varphi^\uparrow(s,h)$, and the left side is
$f^\sharp_s(h\oplus' k)=f^\sharp_s(k\circ' h)=\varphi^\uparrow(s,k\circ' h)$.
So (M2) reads $\varphi^\uparrow(s,k\circ' h)=\varphi^\uparrow(s',k)\circ
\varphi^\uparrow(s,h)$ with $s'=\cod\varphi^\uparrow(s,h)$, exactly (C2).

The three conditions are mutually equivalent under an identity map on data, so
the assignment is an involution-respecting bijection on hom-sets.
\end{proof}

\begin{lemma}[Functoriality]\label{lem:fun}
$\mathcal C(-)$ together with the assignment of Lemma~\ref{lem:hom} is a
functor $\Phi\colon\DCont\to\Cof$.
\end{lemma}

\begin{proof}
\emph{Identities.} The identity container morphism is $(\id_S,\id)$, which maps
to $(\id_{\ob\mathcal C},\id^\uparrow)$, the identity cofunctor.

\emph{Composition.} Let $(f,f^\sharp)\colon D\to D'$ and
$(g,g^\sharp)\colon D'\to D''$. Their container composite has object map
$g\circ f$ and position map
$(g\circ f)^\sharp_s = f^\sharp_s\circ g^\sharp_{fs}\colon
P''(gfs)\to P'(fs)\to P(s)$
(positions pull back through both stages). Under $\Phi$ this is the cofunctor
with object map $g_0\circ f_0$ and lift
$m\mapsto f^\sharp_s\big(g^\sharp_{fs}(m)\big)
=\varphi^\uparrow\big(s,\psi^\uparrow(\varphi_0 s,m)\big)$,
which is exactly the composite cofunctor $\psi\varphi$ of
Definition~\ref{def:cof}. Hence $\Phi$ preserves composition.
\end{proof}

\begin{theorem}[Directed containers are categories-and-cofunctors]\label{thm:main}
The functor $\Phi\colon\DCont\to\Cof$ is an isomorphism of categories:
\[
  \boxed{\ \DCont\ \cong\ \Cof\ }
\]
\end{theorem}

\begin{proof}
$\Phi$ is fully faithful by Lemma~\ref{lem:hom} (bijection on every hom-set)
and a functor by Lemma~\ref{lem:fun}. On objects, $D\mapsto\mathcal C(D)$ is a
bijection: by Theorem~4.1 of \texttt{directed-categories.tex} every small
category is $\mathcal C(D)$ for a unique directed container $D$. A
bijective-on-objects, fully faithful functor is an isomorphism of categories
(its object bijection and hom bijections assemble into a strict inverse).
\end{proof}

\section{Functors live elsewhere: the opposite variance}

Theorem~\ref{thm:main} says the morphisms native to directed containers are
cofunctors, not functors. We make the contrast precise and exhibit where
functors do sit.

\begin{proposition}[The dictionary does not extend to $\Cat$]\label{prop:neq}
Let $\Phi_0\colon\ob\DCont\to\ob\Cat$, $D\mapsto\mathcal C(D)$, be the
object-level dictionary (a bijection on objects, Theorem~4.1). There is
\emph{no} functor $\DCont\to\Cat$ that agrees with $\Phi_0$ on objects and is
full. In particular the object dictionary does not extend to an equivalence
(indeed not to any isomorphism) $\DCont\simeq\Cat$.
\end{proposition}

\begin{proof}
Let $D_{\mathbf1}, D_{\mathbf2}$ be the directed containers of the terminal
category $\mathbf 1$ and $\mathbf 2=\{0\xrightarrow{u}1\}$. By
Theorem~\ref{thm:main} (via Lemma~\ref{lem:hom}),
$|\DCont(D_{\mathbf1},D_{\mathbf2})|=|\Cof(\mathbf1,\mathbf2)|$. We count the
latter. A cofunctor $\mathbf1\co\mathbf2$ picks $\varphi_0(\ast)\in\{0,1\}$ and
must lift \emph{every} outgoing morphism of $\varphi_0(\ast)$ into the unique
morphism $\id_\ast$ out of $\ast$, with (C0) forcing
$\varphi_0(\cod(\text{lift}))=\cod$. Object $0$ has the outgoing $u\colon0\to1$
whose lift would need codomain $\mapsto 1\neq\varphi_0(\ast)$: impossible.
Object $1$ has only $\id_1$, lifting to $\id_\ast$. So
$|\Cof(\mathbf1,\mathbf2)|=1$, whence $|\DCont(D_{\mathbf1},D_{\mathbf2})|=1$.
But $|\Cat(\mathbf1,\mathbf2)|=2$ (choose either object). A functor
$G\colon\DCont\to\Cat$ with $G=\Phi_0$ on objects would restrict to a map
$\DCont(D_{\mathbf1},D_{\mathbf2})\to\Cat(\mathbf1,\mathbf2)$ from a
$1$-element set to a $2$-element set, which cannot be surjective; so $G$ is
not full. As an equivalence (or isomorphism) is in particular full, no such
extension exists.
\end{proof}

\begin{remark}[$\Cof$ versus $\Cat$ as abstract categories]
Proposition~\ref{prop:neq} is the operative fact for the program: the chain's
object equivalence does \emph{not} carry functors. One can ask the sharper
question whether $\Cof$ and $\Cat$ are equivalent as abstract categories
(ignoring the dictionary). The hom-cardinality data already shows they are not
\emph{isomorphic} over their common object class, and moreover the comparison
is not even monotone: for the $3$-chain
$|\Cat(\mathbf3,\mathbf2)|=4<5=|\Cof(\mathbf3,\mathbf2)|$ while
$|\Cat(\mathbf2,\mathbf2)|=3>2=|\Cof(\mathbf2,\mathbf2)|$
(Table~\ref{tab:counts}), so neither sits inside the other. A full
equivalence-level separation requires a relabeling-invariant (e.g.\ comparing
the adjoint profiles of $\Hom(\mathbf1,-)$, which is $\ob$ on $\Cat$ but the
``sink-objects'' functor on $\Cof$); we do not pursue it here, as
Proposition~\ref{prop:neq} already settles the question the dictionary poses.
\end{remark}

\begin{proposition}[Functors = covariant-position morphisms]\label{prop:cov}
Let an \emph{op-variance container morphism} $D\to D'$ be $(f,\phi)$ with
$f\colon S\to S'$ and a \emph{covariant} position map
$\phi_s\colon P(s)\to P'(fs)$ satisfying
\textup{(N0)} $f(s\downarrow p)=(fs)\downarrow'\phi_s(p)$,
\textup{(N1)} $\phi_s(\mathsf o(s))=\mathsf o'(fs)$, and
\textup{(N2)} $\phi_s(p\oplus q)=\phi_s(p)\oplus'\phi_{s\downarrow p}(q)$.
Then $(f,\phi)\mapsto F$, $F=f$ on objects, $F(p)=\phi_s(p)$ on a morphism
$p$ out of $s$, is a bijection onto functors $\mathcal C(D)\to\mathcal C(D')$.
\end{proposition}

\begin{proof}
A functor $F$ is precisely: an object map $f=F$; for each $s$ a map of
out-sets $\phi_s=F|_{P(s)}\colon P(s)\to P'(fs)$; with $F$ preserving codomains
(N0: $F(\cod p)=\cod F(p)$), identities (N1: $F(\id_s)=\id_{Fs}$), and
composition (N2: $F(q\circ p)=F(q)\circ F(p)$, i.e.\
$F(p\oplus q)=\phi_s(p)\oplus'\phi_{s\downarrow p}(q)$). Each datum and law
matches under an identity map on data, giving the bijection.
\end{proof}

So functors and cofunctors are the \emph{two opposite variances} of a
morphism of directed structure: covariant on positions (push forward,
$P(s)\to P'(fs)$) gives functors; contravariant on positions (pull back,
$P'(fs)\to P(s)$, the genuine container morphism) gives cofunctors. Only the
latter is a morphism of containers.

\section{The \texorpdfstring{$\mathbf{Poly}$}{Poly} perspective}\label{sec:poly}

Under the embedding of directed containers as comonoids in
$(\mathbf{Poly},\triangleleft)$ --- the polynomial-comonad incarnation of a
small category --- container morphisms are exactly comonoid morphisms.
Theorem~\ref{thm:main} is then the statement that comonoid morphisms in
$(\mathbf{Poly},\triangleleft)$ between category-comonoids are cofunctors
(Ahman--Uustalu; Garner; Spivak--Garner--Fairbanks). Functors, by
Proposition~\ref{prop:cov}, are \emph{not} comonoid morphisms; they are a
different structure (carried by bicomodules / parametric right adjoints in
$\mathbf{Poly}$). The variance mismatch flagged at the outset is thus not an
accident of bookkeeping but the shadow of two different universal properties:
the comonoid maps (cofunctors) and the ``functor'' maps occupy different
hom-objects of $\mathbf{Poly}$.

\paragraph{Why this matters for the grant.} The equivalence chain advertised in
the seed --- containers $\simeq$ directed containers $\simeq$ polynomial
comonoids $\simeq$ small categories --- is an equivalence \emph{of objects}.
At the level of morphisms it is an isomorphism $\DCont\cong\Cof$, \emph{not}
$\DCont\cong\Cat$. Any application that transports a construction along the
chain (supply-chain composition, agent orchestration, the internal replacement
theorem) must use cofunctors as the maps. Mistaking them for functors silently
changes the morphisms --- and, as Proposition~\ref{prop:neq} shows, even their
\emph{number}. This is precisely the kind of compositional-correctness pitfall
the grant claims to eliminate by formalisation.

\section{Computational verification}

The Python script \texttt{dcont\_morphisms\_check.py} represents a small
category concretely and enumerates, between two categories $C,D$: (A) functors
$C\to D$; (B) cofunctors $C\co D$ (Definition~\ref{def:cof}); (C)
directed-container morphisms via (M0)--(M2) computed \emph{directly} from the
container data $(\mathsf o,\downarrow,\oplus)$ (Definition~\ref{def:dcm}); and
(D) op-variance morphisms via (N0)--(N2). Over all $36$ ordered pairs from
$\{\mathbf1,\mathbf2,\mathbf3,\mathbb Z/2,\mathbb Z/3,M\}$ (with $M$ the
two-element idempotent monoid $\{1,t\},\,t^2=t$):
\[
  \#(C)=\#(B)\ \text{in all }36\text{ cases},\qquad
  \#(D)=\#(A)\ \text{in all }36\text{ cases},
\]
confirming Lemma~\ref{lem:hom} and Proposition~\ref{prop:cov}; while
$\#(A)\neq\#(B)$ in $17$ of the $36$ cases, confirming
Proposition~\ref{prop:neq}. A representative excerpt:

\begin{table}[h]\centering
\begin{tabular}{lrrr}
\toprule
$C\to D$ & \#functors & \#cofunctors & \#DCont-morphisms\\
\midrule
$\mathbf1\to\mathbf2$ & $2$ & $1$ & $1$\\
$\mathbf2\to\mathbf2$ & $3$ & $2$ & $2$\\
$\mathbf2\to\mathbb Z/3$ & $3$ & $1$ & $1$\\
$\mathbf3\to\mathbf2$ & $4$ & $5$ & $5$\\
$\mathbf3\to\mathbf3$ & $10$ & $6$ & $6$\\
$\mathbb Z/3\to\mathbb Z/3$ & $3$ & $3$ & $3$\\
$\mathbf3\to M$ & $4$ & $5$ & $5$\\
\bottomrule
\end{tabular}
\caption{Functors vs.\ cofunctors vs.\ directed-container morphisms.
\#cofunctors $=$ \#DCont-morphisms throughout; \#functors differs.
$\mathbf3\to\mathbf2$ and $\mathbf2\to\mathbf2$ show neither of
$\Cat,\Cof$ dominates the other.}
\label{tab:counts}
\end{table}

\section{Conclusion}

The dictionary is now complete at the morphism level. Objects: directed
containers $=$ small categories. Morphisms: container morphisms $=$
\emph{cofunctors}. The assignment is an \emph{isomorphism of categories}
$\DCont\cong\Cof$. The tempting $\DCont\simeq\Cat$ is false; functors are the
opposite variance and are not morphisms of containers. The next step is to
formalise $\Cof$ and Theorem~\ref{thm:main} in Lean (M2/M4), and to revisit
every place in the applied notes where ``functor between directed containers''
was written, replacing it with the correct cofunctor.

\end{document}
